\documentclass[11pt,a4paper]{article}

\usepackage{fontspec}
\usepackage{polyglossia}
\setdefaultlanguage{russian}
\setotherlanguage{english}
\setmainfont{DejaVu Serif}
\setsansfont{DejaVu Sans}
\setmonofont{DejaVu Sans Mono}

\usepackage[a4paper,left=2cm,right=2cm,top=1.8cm,bottom=1.8cm]{geometry}
\usepackage{setspace}
\usepackage{enumitem}
\usepackage{graphicx}
\usepackage[hidelinks]{hyperref}
\usepackage{xurl}
\usepackage{float}

\setstretch{1.08}
\setlength{\parindent}{1.25cm}
\setlength{\parskip}{0.25em}
\setlist{itemsep=0.3em,topsep=0.4em}
\pagestyle{plain}

\begin{document}

\begin{center}
    {\Large\bfseries Отчёт по лабораторной работе по созданию агентов}

    \vspace{0.4em}
    {\large по дисциплине «Технология разработки программного обеспечения»}

    \vspace{0.7em}
    {\Large\bfseries «Агент для поиска первоисточника новостей»}
\end{center}

\section*{Исполнители}

\begin{itemize}[leftmargin=1.8em]
    \item Джабраилов Арсен,
    \item Ибатуллин Денис,
    \item Пешехонова Ирина,
    \item Рузанова Анастасия,
    \item Смирнов Никита,
    \item Юркевич Анастасия.
\end{itemize}

\section*{Ссылки на проект}

\begin{itemize}[leftmargin=1.8em]
    \item \textbf{Репозиторий проекта:} \url{https://github.com/Irene-Pesh-code/NotFakeNews-agent}
    \item \textbf{Telegram-бот:} \url{https://t.me/not_fake_news_agent_bot}
\end{itemize}

\section*{Цель и назначение проекта}

Цель проекта --- создать агента, который помогает находить первоисточник новости. Пользователь отправляет боту публикацию, после чего агент ищет более ранние материалы с тем же содержанием, переходит по ссылкам на упомянутые источники и сравнивает даты публикаций. В результате пользователь получает ссылку на найденный первоисточник и может посмотреть, откуда началось распространение новости.

Такой агент нужен, потому что одна и та же новость часто появляется сразу в нескольких Telegram-каналах и интернет-изданиях. Во время перепечатки ссылка на исходную публикацию может потеряться, а текст может быть сокращён или пересказан другими словами. Искать источник вручную долго: требуется открыть несколько сайтов, сравнить формулировки, проверить даты и пройти по цепочке ссылок. Агент выполняет эти действия автоматически и собирает результат в одном сообщении.

Проект также позволяет увидеть, как новость распространялась между источниками. Помимо первоисточника, бот показывает другие подтверждённые публикации, их даты и фрагменты текста. Пользователь получает не только готовую ссылку, но и понятное объяснение, на основании каких материалов был сформирован ответ.

\section*{Что принимает агент}

Для запуска поиска пользователю достаточно отправить боту один из следующих вариантов:

\begin{itemize}[leftmargin=1.8em]
    \item пересланный пост из Telegram-канала;
    \item ссылку на публичный пост Telegram;
    \item ссылку на новостную статью или другую веб-страницу;
    \item текст новости;
    \item отдельную цитату, по которой требуется найти публикацию.
\end{itemize}

Для пересланного поста агент получает текст, название канала и дату исходной публикации из данных Telegram. По ссылке на сайт или публичный пост он загружает страницу и извлекает её содержимое. Если пользователь отправляет обычный текст, поиск начинается по ключевым цитатам, именам, организациям, датам и числам из этого текста.

\section*{Как работает агент}

Работа агента состоит из нескольких последовательных этапов.

\begin{enumerate}[leftmargin=2em]
    \item \textbf{Разбор сообщения.} Агент определяет тип входных данных, извлекает текст, дату, название источника и ссылки, которые находятся внутри публикации.

    \item \textbf{Выделение ключевых фрагментов.} Языковая модель выбирает от одного до трёх наиболее полезных фрагментов новости. Обычно это прямая цитата или утверждение с именами, названиями организаций, датами и числами. Также формируются поисковые запросы.

    \item \textbf{Поиск ссылок на другие источники.} В тексте находятся выражения «как сообщает», «по данным», «со ссылкой на», «передаёт» и похожие конструкции. Агент связывает их с URL или названием указанного источника.

    \item \textbf{Переход по цепочке публикаций.} Агент открывает найденные ссылки и повторяет анализ для каждой новой страницы. Так строится цепочка от более поздних публикаций к более ранним.

    \item \textbf{Дополнительный веб-поиск.} Если в тексте указан источник без ссылки или цепочка публикаций заканчивается, агент выполняет поиск по цитате, названию источника и ключевым фактам новости.

    \item \textbf{Проверка найденных страниц.} Для каждой открытой страницы программа проверяет наличие цитаты или совпадение ключевых фактов, определяет дату публикации и сохраняет подходящий фрагмент текста.

    \item \textbf{Выбор первоисточника.} Публикации сравниваются по датам и типу совпадения. Первоисточником выбирается самая ранняя подтверждённая публикация.

    \item \textbf{Формирование ответа.} Бот показывает первоисточник, дату, подтверждающий фрагмент, дополнительные публикации, статистику поиска и уровень уверенности.
\end{enumerate}

\section*{Инструменты агента}

Для управления поиском языковой модели доступны четыре инструмента:

\begin{itemize}[leftmargin=1.8em]
    \item \texttt{search} --- выполняет веб-поиск по заданному запросу и возвращает список найденных материалов;
    \item \texttt{open} --- открывает выбранный результат поиска, извлекает текст и дату, проверяет совпадение новости и находит ссылки на другие источники;
    \item \texttt{follow} --- переходит по ссылке, указанной в уже открытой публикации, и добавляет новую страницу в цепочку;
    \item \texttt{finish} --- завершает поиск и передаёт идентификатор выбранного первоисточника для формирования результата.
\end{itemize}

Модель самостоятельно выбирает последовательность вызова этих инструментов. Например, она может сначала перейти по ссылке из исходной публикации, затем выполнить поиск по точной цитате и параллельно открыть несколько подходящих результатов. После каждого действия модель получает сведения о найденной странице и решает, какой шаг выполнить дальше.

\section*{Проверка найденных публикаций}

Сначала программа ищет выбранную цитату в тексте страницы. Перед сравнением текст нормализуется. Для немного изменённых формулировок используется нечёткое сравнение. Числа при этом проверяются отдельно, поэтому публикации с другими датами, процентами или количественными значениями не считаются одинаковыми.

Новость также может быть пересказана другими словами. Для таких случаев агент сравнивает ключевые факты: числа, имена собственные и значимые слова. Найденный фрагмент сохраняется и затем показывается пользователю. По типу проверки совпадение отмечается как точное, близкое или фактологическое.

\section*{Ограничения агента}

Чтобы поиск завершался за разумное время и не выполнял слишком много внешних запросов, для одного запуска установлены следующие ограничения:

\begin{itemize}[leftmargin=1.8em]
    \item не более 16 шагов языковой модели;
    \item не более 6 запросов к веб-поиску;
    \item не более 25 загруженных страниц;
    \item не более 6 переходов в глубину по цепочке источников;
    \item до 6 результатов на один поисковый запрос;
    \item около 220 секунд на полный запуск агента при стандартных настройках;
    \item размер загружаемой HTML-страницы --- до 3 Мбайт;
    \item до 10 дополнительных источников в итоговом сообщении Telegram.
\end{itemize}

Для одного пользователя одновременно выполняется одна проверка. Пока агент обрабатывает сообщение, бот показывает текущий этап работы и предлагает дождаться завершения поиска перед отправкой следующей новости.

\section*{Как использовать бота}

\begin{enumerate}[leftmargin=2em]
    \item Открыть Telegram-бота.
    \item Нажать кнопку запуска или отправить команду \texttt{/start}. Бот покажет краткую справку о поддерживаемых форматах.
    \item Переслать пост, отправить ссылку или вставить текст новости одним сообщением.
    \item Дождаться завершения поиска. Обычно обработка занимает от нескольких секунд до нескольких минут. Текущий этап отображается в сообщении со статусом.
    \item Изучить результат: открыть первоисточник, посмотреть дату и сравнить найденный фрагмент с исходной новостью.
\end{enumerate}

В итоговом сообщении указываются найденный первоисточник, дата публикации, фрагмент текста, количество проверенных страниц, список дополнительных источников и уровень уверенности. Высокий уровень соответствует точному совпадению и определённой дате. Средний уровень используется для близкого совпадения, пересказа или даты, известной только с точностью до дня. Если поиск не определил первоисточник, бот показывает публикацию, до которой удалось дойти, и низкий уровень уверенности.

\section*{Пример работы}

На рисунке~\ref{fig:bot-example} изображен скриншот диалога с ботом: сообщение пользователя и полученный результат поиска.

\begin{figure}[H]
    \centering
    \includegraphics[width=0.97\textwidth]{images/example_yandex.png}
    \caption{Пример использования агента для поиска первоисточника новости}
    \label{fig:bot-example}
\end{figure}

\section*{Использованные технологии}

Проект написан на Python~3.11. Telegram-интерфейс реализован с помощью библиотеки \texttt{aiogram}, HTTP-запросы выполняются через \texttt{httpx}, а основной текст страниц извлекается библиотекой Trafilatura. Для управления поиском используется языковая модель через OpenAI-совместимый API, для веб-поиска --- сервис Tavily, для сравнения текста --- RapidFuzz. Проект можно запускать локально или в Docker Compose. Для проверки отдельных компонентов предусмотрен консольный скрипт, а алгоритмы сопоставления цитат, фактов и URL покрыты автоматическими тестами.

\end{document}
